% body.tex — P03: Zigzag Persistence, Survey Design, and Labour Market Fragmentation
% Converted from drafts/v2-2026-03.md

%%%%%%%%%%%%%%%%%%%%%%%%%%%%%%%%%%%%%%%%%%%%%%%%%%%
\section{Introduction}
%%%%%%%%%%%%%%%%%%%%%%%%%%%%%%%%%%%%%%%%%%%%%%%%%%%

\subsection{Motivation}

Labour markets restructure during economic crises.
The 2008 Global Financial Crisis displaced over 750,000 UK workers into
unemployment within eighteen months, shifted sectoral employment towards
low-paid services, and inaugurated a decade of real-wage stagnation without
post-war precedent \citep{gregg2010unemployment, blundell2014wages}.
Cross-sectional measures (unemployment rates, Gini coefficients, median wages)
capture the level of such shocks.
Longitudinal sequence analysis captures their effect on individual careers
\citep{abbott1995sequence, halpin1998class}.
Neither captures whether crises \emph{restructure the space of possible
trajectories}: whether the set of career types available to individuals is
geometrically different after a crisis than before.

Persistent homology (PH) provides a multi-scale summary of the topology of a
data set (its connected components, loops, and higher-dimensional voids) that is
stable under perturbation and robust to noise
\citep{edelsbrunner2010computational, carlsson2009topology}.
In a companion paper \citep{dorman2026a}, we apply standard PH to a static
cross-section of UK employment-income trajectories and identify significant
topological structure exceeding first-order Markov baselines under Wasserstein
testing.

But a single persistence diagram cannot distinguish structural change from stable
heterogeneity.
A trajectory space that looks topologically complex in 2022 might have been
identical in 1991.
To answer temporal questions, we need methods that track topology across time.

\subsection{Zigzag Persistence}

Zigzag persistence \citep{carlsson2010zigzag} generalises PH to sequences of
spaces connected by inclusions in alternating directions:
\[
  K_1 \hookrightarrow K_1 \cup K_2 \hookleftarrow K_2
       \hookrightarrow K_2 \cup K_3 \hookleftarrow K_3
       \hookrightarrow \cdots
\]
The zigzag algorithm computes a single barcode over this tower.
A bar of length $k$ in the $\HH{0}$ barcode corresponds to a connected component
that remained distinct across $k$ consecutive snapshots---qualitatively different
from per-year PH.

Zigzag persistence has been applied in neuroscience
\citep{bendich2016persistent}, dynamic networks \citep{kim2021analysis}, and
materials science \citep{townsend2020representation}.
Related sliding-window methods have been used in financial time series:
\citet{gidea2018topological} detect early warnings of the 2008 crash via
persistence landscape norms on rolling windows of stock returns.
To our knowledge, zigzag persistence has not been applied to longitudinal
socioeconomic panel data.

\subsection{The Identification Problem}

When longitudinal data spans multiple survey instruments, as when the British
Household Panel Study (BHPS, 1991--2008) gives way to Understanding Society
(USoc, 2009 onwards), any change in topology could reflect genuine structural
change, survey-design differences (sampling frame, questionnaire wording, income
measurement), or both.
Every time-series analysis of linked panels faces this identification problem
\citep{lynn2009methodology}, but topological methods make it unusually visible:
the Wasserstein distance between annual persistence diagrams quantifies how
different two years look in topology space, and this distance responds to both
substantive and design-related variation.

We exploit this property, treating the survey transition as a natural experiment.
If post-2008 topological change is economic (the GFC), it should affect spanning
individuals (present in both surveys) and newcomers alike.
If it is survey-design (the USoc sampling frame captures a broader population),
spanning individuals should maintain pre-2008 topology while newcomers drive the
change.
As we show, the data decisively support the second interpretation.

\subsection{Research Questions}

\begin{enumerate}
  \item \textbf{Temporal evolution.} How does the topology of the UK trajectory
    space change across 32 annual snapshots (1991--2022)?
  \item \textbf{Source identification.} Is the dominant post-2008 topological
    change driven by economic restructuring or by the BHPS-to-USoc survey
    transition?
  \item \textbf{Within-era dynamics.} Within a single survey (BHPS 1991--2008),
    does trajectory topology correlate with macroeconomic conditions?
\end{enumerate}

\subsection{Contribution}

We make three contributions.
First, we apply zigzag persistence to population-level panel data and show that
its headline finding (a topological discontinuity at 2008) is largely an artefact
of survey design, not the Global Financial Crisis.
Zigzag persistence works, but it detects \emph{all} sources of distributional
change, including non-substantive ones.
Second, we develop a set of diagnostics for telling apart genuine structural
change and survey confounds in time-varying TDA: spanning-individual
decomposition, era-specific zigzag, sub-sampled Betti comparison, and pool-draw
null models.
Third, having controlled for survey design, we document a real secular trend in
topological complexity during 1991--2008 that tracks labour market composition
(employment rate, part-time share, self-employment) rather than business-cycle
dynamics (GDP growth, inflation).

Section~\ref{sec:lit} reviews the relevant literature.
Section~\ref{sec:data} describes the data, embedding, and analytical methods.
Section~\ref{sec:results} presents results: the na\"ive zigzag findings
(\S\S\ref{sec:naive}--\ref{sec:sensitivity}), the survey-confound decomposition
(\S\ref{sec:decomp}), the Wasserstein distance analysis with null models
(\S\ref{sec:wasserstein}), and the macroeconomic correlations
(\S\ref{sec:macro}).
Section~\ref{sec:discussion} discusses implications.

%%%%%%%%%%%%%%%%%%%%%%%%%%%%%%%%%%%%%%%%%%%%%%%%%%%
\section{Related Work}
\label{sec:lit}
%%%%%%%%%%%%%%%%%%%%%%%%%%%%%%%%%%%%%%%%%%%%%%%%%%%

\subsection{Business Cycles and Labour Market Structure}

Economic recessions reshape not just unemployment levels but the structure of
labour market transitions.
\citet{elsby2010labor} document that the 2008 recession produced both elevated
job-loss rates and reduced job-finding rates, with the latter persisting well
beyond the trough.
In the UK, \citet{gregg2010unemployment} show the GFC disproportionately
affected younger workers and manufacturing, while \citet{blundell2014wages}
document falling real wages alongside rising employment.
\citet{bukodi2019social} use sequence analysis on British birth cohorts to show
declining intergenerational class fluidity, and \citet{halpin1998class} identify
five trajectory types with distinct transition patterns in early BHPS data.

These approaches share an individual-level or typological focus.
None can detect whether the \emph{global geometry} of the trajectory space
changes during crises.

\subsection{Time-Varying Topology}

Extensions of PH to temporal data fall into three categories.
\textbf{Sliding-window PH} \citep{perea2015sliding, perea2019topological}
embeds a time series via delay coordinates and computes PH per window.
\citet{gidea2018topological} apply this to daily stock returns, detecting
pre-crash signals in persistence landscape norms.
This operates on rolling windows of a single signal; ours operates on annual
population cross-sections.

\textbf{Zigzag persistence} \citep{carlsson2010zigzag} computes a barcode over
a tower of simplicial complexes connected by alternating inclusions.
The diamond construction (Cohen-Steiner et al.\ 2009) is standard: given
complexes $K_1, \ldots, K_T$, form $K_1 \hookrightarrow K_1 \cup K_2
\hookleftarrow K_2 \hookrightarrow \cdots$.
Applications include brain-state tracking \citep{bendich2016persistent},
evolving random graphs \citep{dey2019simba}, and dynamic molecular structures
\citep{townsend2020representation}.

\textbf{Vineyard persistence} \citep{cohen2006vines} tracks pairs as a
filtration varies continuously, less natural for discrete population snapshots
with different individuals.

\subsection{TDA in the Social Sciences}

Social-science TDA applications remain sparse.
\citet{feng2021persistent} use PH on US Congressional voting networks to detect
partisan polarisation.
\citet{iacopini2019simplicial} use simplicial complexes for social contagion.
These are cross-sectional network studies; ours uses longitudinal individual-level
trajectories with zigzag tracking.

\subsection{Survey Design and Longitudinal Comparability}

The transition from BHPS to USoc introduced changes in sampling frame (USoc
covers a larger population with ethnic minority boost samples), questionnaire
design, interview mode (more online, less face-to-face), and income measurement
\citep{lynn2009methodology, buck2011understanding}.
Studies linking the two surveys routinely find discontinuities attributable to
design rather than substance \citep{benzeval2020understanding}.
Persistent homology provides an independent, geometric measure of this
discontinuity.

%%%%%%%%%%%%%%%%%%%%%%%%%%%%%%%%%%%%%%%%%%%%%%%%%%%
\section{Data and Methods}
\label{sec:data}
%%%%%%%%%%%%%%%%%%%%%%%%%%%%%%%%%%%%%%%%%%%%%%%%%%%

\subsection{Data}

We draw on the British Household Panel Study (BHPS, 1991--2008) and
Understanding Society (USoc, 2009--2022), the longest continuous UK household
panel.
The sample comprises all individuals observed for at least 10 consecutive years,
yielding $N = 35{,}789$ trajectories: 8,509 from the BHPS era and 27,280 from
USoc.
Of these, 8,459 ``spanning'' individuals appear in both surveys.

Each trajectory is a time-ordered sequence of annual states from a nine-element
state space combining employment status (employed, unemployed, inactive) with
income band (low, mid, high), where ``low'' is below 60\% of the contemporary
median equivalised household income.

Annual active counts range from 4,729 (1991) to 28,632 (2018).
The BHPS--USoc transition produces a near-threefold jump: from 9,587 (2008) to
27,175 (2009).
This discontinuity is central to the paper.

\subsection{Embedding}

Following \citet{dorman2026a}, each trajectory is encoded as an $n$-gram
frequency vector (counts of all state unigrams, bigrams, and trigrams from the
nine-state alphabet, dimension $9 + 81 + 729 = 819$) then projected into
$\mathbb{R}^{20}$ via PCA.
The PCA loadings are \textbf{frozen}: fitted once to the full sample and applied
to all annual sub-selections without re-fitting.
This ensures a common coordinate system across years.

\subsection{Independent Annual Persistent Homology}

For each year $Y \in [1991, 2022]$, we select individuals active in that year
from the embedded point cloud.
To control for sample-size variation, we sub-sample to $n = 8{,}000$ per year
(with 10--15 repetitions) and select $L = 300$ maxmin landmarks per sub-sample.
We compute Vietoris--Rips PH via GUDHI \citep{maria2014gudhi} with max edge
length 1.1.

From each annual persistence diagram, we extract:
\begin{itemize}
  \item $\AUC(\beta_0)$: area under the Betti-0 curve integrated over
    $\varepsilon \in [0.3, 0.8]$
  \item $\TP(k)$: sum of bar lifetimes (total persistence) in dimension $k$
  \item $\PE(k)$: Shannon entropy of the normalised lifetime distribution in
    dimension $k$
\end{itemize}

\subsection{Zigzag Persistence}

We compute zigzag persistence via the diamond construction over 32 annual
sub-clouds.
For each year, $L$ points are selected by maxmin greedy subsampling.
Unique landmarks across all years form the vertex set.
A vertex representing individual $I$ (active in years $[s_I, e_I]$) is present
at zigzag levels spanning that interval; edges between individuals $I, J$ with
$d(v_I, v_J) \leq \varepsilon$ are present for the intersection of their
vertex-level intervals.
We use dionysus~2 \citep{morozov2010dionysus} with default parameters
$\varepsilon = 2.0$, sparse $= 0.5$, max dimension 1.

To assess sensitivity, we compute zigzag over a 2D grid: 11 values of
$\varepsilon$ (0.34--1.00) $\times$ 4 values of $L$ (75, 100, 150, 200), plus
era-specific runs at matched sample sizes.

\subsection{Wasserstein Distance Matrix}

We compute pairwise 2-Wasserstein distances between all 32 annual $\HH{0}$
persistence diagrams using GUDHI:
\[
  \Wass(D_Y, D_{Y'}) \quad \text{for all } Y, Y' \in [1991, 2022].
\]
Both full-sample (300 landmarks) and sub-sampled ($n = 8{,}000$, $L = 300$, 10
repetitions, mean distances) versions are computed.
We analyse the resulting $32 \times 32$ matrix via hierarchical clustering
(Ward's method), Mantel tests (correlation between $\Wass$ and temporal distance
$|Y - Y'|$), and multidimensional scaling.

\subsection{Null Models}

We employ three null models.

\textbf{Label permutation} ($n = 1{,}000$).
Permute the year labels on the observed $\Wass$ matrix, keeping the first 18
positions as ``BHPS'' and the last 14 as ``USoc.''
Compute block ratio (cross-era/within-era mean $\Wass$), consecutive-year mean,
and Mantel $r$ for each permutation.
Tests whether temporal structure is significant.

\textbf{Pool-draw pairs} ($n = 100$).
Draw two random subsets of $n = 8{,}000$ from the full pool (ignoring year
labels), compute fresh Rips diagrams, and compute $\Wass$.
Tests baseline diagram variability under year-ignorant sampling.

\textbf{Pool-draw full} ($n = 50$, 32 diagrams each).
Draw 32 independent random subsets, assign them to ``BHPS'' (first 18) and
``USoc'' (last 14) positions, and compute block statistics.
Tests whether year-specific population composition produces block structure.

\subsection{Spanning-Individual Decomposition}

To distinguish survey-coverage effects from economic change, we partition the
USoc-era sample into:
\begin{itemize}
  \item \textbf{Spanning individuals} ($n = 8{,}459$): present in both BHPS and
    USoc
  \item \textbf{USoc-only newcomers} ($n = 21{,}551$): present only in USoc
\end{itemize}
We compute sub-sampled Betti curves and persistence summaries separately for each
group and compare them to the BHPS reference distribution.
If the post-2008 topological shift is economic, spanning individuals should show
it.
If it is survey coverage, only newcomers should differ from BHPS.

\subsection{Macroeconomic Correlations}

We correlate six topological time series with seven UK macroeconomic indicators
(1991--2022) from ONS: real GDP growth, ILO unemployment rate, Gini coefficient,
employment rate (16--64), part-time employment share, self-employment share, and
CPI inflation.

Four specifications control for confounds:
\begin{enumerate}
  \item \textbf{Raw Pearson correlations} (full panel, 32 years)
  \item \textbf{Within-era} (BHPS-only, USoc-only)---eliminates
    survey-transition effects
  \item \textbf{Detrended} (residuals from linear time regression)---removes
    shared secular trends
  \item \textbf{First-differenced} ($\Delta Y \sim \Delta X$)---most
    conservative, removes all level effects
\end{enumerate}
All $\pval$-values are corrected for multiple testing via Benjamini--Hochberg
FDR at $\alpha = 0.05$.

%%%%%%%%%%%%%%%%%%%%%%%%%%%%%%%%%%%%%%%%%%%%%%%%%%%
\section{Results}
\label{sec:results}
%%%%%%%%%%%%%%%%%%%%%%%%%%%%%%%%%%%%%%%%%%%%%%%%%%%

\subsection{The Na\"{i}ve Zigzag: An Apparently Compelling Story}
\label{sec:naive}

At 150 maxmin landmarks and $\varepsilon = 1.0$, the zigzag barcode contains four
$\HH{0}$ bars (Table~\ref{tab:zigzag_bars}).

\begin{table}[htbp]
  \centering
  \caption{Zigzag $\HH{0}$ bars (150 landmarks, $\varepsilon = 1.0$).}
  \label{tab:zigzag_bars}
  \begin{tabular}{clcccp{5cm}}
    \toprule
    Bar & Birth & Death & Lifetime & Interpretation \\
    \midrule
    1 & 1991   & 2022.5 & 31.5 yr & Dominant component \\
    2 & 1991   & 2003   & 12.0 yr & Pre-2000s cluster \\
    3 & 1991   & 1996.5 &  5.5 yr & ERM recession fragment \\
    4 & 2008.5 & 2022.5 & 14.0 yr & GFC-born cluster \\
    \bottomrule
  \end{tabular}
\end{table}

Bar~1 spans the full panel: one connected component persists throughout.
Bar~3 is born with the panel and dies during the mid-1990s recovery, plausibly
the topological signature of the 1993 recession.
Bar~4 is the most striking: it appears at the onset of the GFC and persists to
the end of the panel, a seemingly permanent topological scar.
Unlike the transient 1993 fragment, the GFC cluster never merges back.

Per-year Betti-0 estimates are flat across all 32 years (range 234--274,
CV $<$ 5\%).
Independent annual PH detects no recession signature, which suggests zigzag
reveals cross-year structure invisible to static snapshots.

This is the story of our earlier analysis \citep{dorman2026b}.
It is wrong.

\begin{figure}[htbp]
  \centering
  \includegraphics[width=\linewidth]{fig1_zigzag_barcode}
  \caption{Zigzag $\HH{0}$ barcode (150 landmarks, $\varepsilon = 1.0$) annotated
    with economic events and the BHPS--USoc survey boundary.
    Horizontal bars show birth--death intervals; coloured bands indicate recession years;
    the dashed vertical line marks the survey transition.}
  \label{fig:barcode}
\end{figure}

\subsection{The Sensitivity Grid: Robustness of Zigzag Features}
\label{sec:sensitivity}

A systematic 2D parameter sweep ($11\ \varepsilon$ values $\times\ 4$ landmark
counts) reveals the dependence structure (Table~\ref{tab:sensitivity}).

\begin{table}[htbp]
  \centering
  \caption{$\HH{0}$ bar count across $\varepsilon \times L$ grid (selected rows).}
  \label{tab:sensitivity}
  \begin{tabular}{lrrrr}
    \toprule
    $\varepsilon \backslash L$ & 75 & 100 & 150 & 200 \\
    \midrule
    0.49 & 416 & 490 & 528 & 519 \\
    0.54 & 329 & 357 & 310 & 275 \\
    0.70 &  68 &  51 &  35 &  26 \\
    1.00 &   5 &   4 &   4 &   4 \\
    \bottomrule
  \end{tabular}
\end{table}

Below the per-year Betti knee ($\varepsilon \approx 0.54$), $\HH{0}$ increases
with $L$: noise dominates.
At the knee, $\HH{0}$ is approximately constant across $L$, the regime of
maximal resolution-stability.
Above the knee, $\HH{0}$ decreases with $L$: signal dominates.
The 4-bar structure at $\varepsilon = 1.0$ is consistent at $L = 100$, 150, and
200 (identical bars), so the coarse topology is stable above the knee.
At every parameter combination, features cluster around $2008 \pm 1$, where both
the GFC and the survey transition fall.

\subsection{Decomposition: Survey Design, Not Economic Shock}
\label{sec:decomp}

Three lines of evidence demonstrate that the dominant post-2008 topological shift
reflects the BHPS--USoc survey transition.

\subsubsection{Era-Specific Zigzag}

Running zigzag separately within each survey era (BHPS-only 1991--2008;
USoc-only 2009--2022) at $\varepsilon = 1.0$, $L = 150$:

\begin{center}
  \begin{tabular}{lrl}
    \toprule
    Era & $\HH{0}$ bars & Structure \\
    \midrule
    BHPS-only & 3 & All full-span (1991--2008.5) \\
    USoc-only & 4 & All full-span (2009--2022.5) \\
    Full panel & 4 & Deaths at 1996.5, 2003; birth at 2008.5 \\
    \bottomrule
  \end{tabular}
\end{center}

No within-era merger events occur: all era-specific bars span their full era.
The temporal features in the full-panel zigzag (bars dying at 1996.5 and 2003;
the GFC birth at 2008.5) arise from cross-era interactions when spanning
respondents bridge the two surveys' distinct subpopulations.
Bar~4, the ``GFC cluster,'' is a product of pooling surveys with different
sampling frames.

Sub-sampling USoc to match BHPS sample sizes ($n = 9{,}500$, 10 random draws)
reduces the component count from 4 to a mean of 2.2 (range 1--3).
The fourth component depends on sample size at coarse scales.

\subsubsection{Spanning-Individual Decomposition}

Spanning individuals ($n = 8{,}459$), present in both surveys, serve as an
internal control.

\begin{center}
  \begin{tabular}{lll}
    \toprule
    Population & $\beta_0(0.70)$, 2009--2015 & vs.\ BHPS reference ($\approx 6.0$) \\
    \midrule
    Spanning only  & $\approx 5.7$ & Matches BHPS \\
    USoc newcomers & $\approx 9.5$ & 50\% more complex \\
    Combined       & $\approx 8.1$ & Intermediate \\
    \bottomrule
  \end{tabular}
\end{center}

Spanning individuals maintain BHPS-like topology after the survey transition.
Their embedding positions are stable: only 4.6\% fall outside the BHPS-era
manifold (vs.\ a 5\% baseline).
USoc-only newcomers ($n = 21{,}551$) look different: 11.9\% fall outside the
BHPS manifold (Mann--Whitney $p < 10^{-6}$), and their PCs 1--4 are shifted.
State-composition analysis shows that USoc newcomers have more Employed-Low
(5.6\% vs 3.4\% among spanners) and more Inactive-High (12.5\% vs 8.5\%)
trajectories, career types that the BHPS sampling frame underrepresented.

\begin{figure}[htbp]
  \centering
  \includegraphics[width=\linewidth]{fig3_spanning_decomposition}
  \caption{Spanning individuals vs.\ USoc newcomers: embedding shift (\% outside
    BHPS-era manifold) and Betti-0 comparison at matched sample sizes.}
  \label{fig:spanning}
\end{figure}

\subsubsection{Sub-sampled Betti Curves}

$\AUC(\beta_0, \varepsilon = 0.3\text{--}0.8)$ at matched $n = 8{,}000$ shows a
gradual BHPS-era rise (8.95 in 1991 $\to$ 11.49 in 2008), a step up at 2009
($\to 14.82$), and a plateau through 2022.
The step at 2009 persists at matched sample sizes.

Total persistence $\HH{0}$ shows a clear BHPS-era trend ($+0.31$/year,
$p < 0.0001$) followed by a flat USoc era ($-0.05$/year, $p = 0.49$).
CUSUM changepoint analysis locates the total-persistence trend break at
2006---two years before the GFC---and the entropy changepoint at 2008.

$\HH{1}$ loops show gradual increase with no discontinuity at the GFC
(pre- vs post-GFC $p = 0.28$).
The loop topology is stable across the transition.

\textbf{Interpretation.}
The post-2008 complexity shift is survey coverage, not economic shock.
Within BHPS alone, a genuine secular trend in topological complexity was already
underway, consistent with the broader narrative of labour market flexibilisation
(growth of part-time work, self-employment, atypical contracts) during the 1990s
and 2000s.
The USoc transition did not create this trend; it brought into the sample a
broader population that the BHPS frame had missed.

\begin{figure}[htbp]
  \centering
  \includegraphics[width=\linewidth]{fig2_topological_timeseries}
  \caption{Sub-sampled topological time series (1991--2022), $n = 8{,}000$ per
    year, 15 repetitions, with BHPS/USoc era shading and CUSUM changepoint
    markers.
    Top: $\AUC(\beta_0)$; middle: total persistence $\TP(\HH{0})$; bottom:
    persistence entropy.}
  \label{fig:timeseries}
\end{figure}

\subsection{Wasserstein Distance Matrix and Null Models}
\label{sec:wasserstein}

\subsubsection{Block Structure}

The sub-sampled $32 \times 32$ $\Wass$ matrix between annual $\HH{0}$ diagrams
shows clean two-block structure (Table~\ref{tab:block_means}).

\begin{table}[htbp]
  \centering
  \caption{Block means of pairwise $\Wass$ (sub-sampled, $\HH{0}$).}
  \label{tab:block_means}
  \begin{tabular}{lrr}
    \toprule
    Block & Mean $\Wass$ & SD \\
    \midrule
    BHPS--BHPS & 0.486 & 0.042 \\
    BHPS--USoc & 0.752 & 0.105 \\
    USoc--USoc & 0.466 & 0.033 \\
    \bottomrule
  \end{tabular}
\end{table}

Cross-era distances are $1.55\times$ within-BHPS and $1.61\times$ within-USoc.
Within-era distances are comparable; the asymmetry is concentrated in the
cross-era block.

Hierarchical clustering (Ward's method) at $k = 2$ perfectly separates BHPS from
USoc: zero misclassifications across 32 years.
At $k = 3$, the algorithm splits BHPS into early (1991--92, 94--96, 99, 2001)
and late (1993, 97--98, 2000, 02--08) sub-eras.
USoc years are not split until $k = 4$ (2009--10 and 2022 separate from the
2011--2021 core).

\begin{figure}[htbp]
  \centering
  \includegraphics[width=0.9\linewidth]{fig4_wasserstein_heatmap}
  \caption{Wasserstein-2 distance heatmap ($32 \times 32$ annual $\HH{0}$
    diagrams, sub-sampled) with Ward's dendrogram.
    BHPS years (1991--2008) are coloured blue; USoc years (2009--2022) red.
    The lower-left BHPS block and upper-right USoc block are visibly cooler
    (lower $\Wass$) than the cross-era blocks.}
  \label{fig:heatmap}
\end{figure}

\subsubsection{Temporal Structure}

Mantel tests reveal strong full-panel time-topology correlation and an
informative within-era asymmetry (Table~\ref{tab:mantel}).

\begin{table}[htbp]
  \centering
  \caption{Mantel test results: correlation of $\Wass$ with temporal distance.}
  \label{tab:mantel}
  \begin{tabular}{lrr}
    \toprule
    Panel & Pearson $r$ & $p$ \\
    \midrule
    Full (1991--2022)  & 0.768 & $< 10^{-6}$ \\
    BHPS (1991--2008)  & 0.423 & $< 10^{-6}$ \\
    USoc (2009--2022)  & 0.144 & 0.17 \\
    \bottomrule
  \end{tabular}
\end{table}

BHPS-era diagrams show moderate temporal ordering: adjacent years are more
similar than distant ones.
USoc-era diagrams are topologically homogeneous: 14 years of persistence diagrams
with no significant temporal structure.
This fits a picture of USoc maintaining a stable expanded sampling frame; once
the broader population is in the sample, the topology stops evolving.

The $2008 \to 2009$ consecutive jump ($\Wass = 0.507$) is the 4th largest of 31
consecutive-year steps, behind $1992 \to 93$ (0.546), $2005 \to 06$ (0.538), and
$1993 \to 94$ (0.515).
It is not an outlier, which fits a survey transition rather than a discrete
economic shock.

\subsubsection{Null Models}

\textbf{Label permutation} ($1{,}000$ permutations).
All six test statistics reject the null (Table~\ref{tab:null_permutation}); the
temporal structure in the $\Wass$ matrix is not random.

\begin{table}[htbp]
  \centering
  \caption{Label-permutation null model ($n = 1{,}000$ permutations).}
  \label{tab:null_permutation}
  \begin{tabular}{lrrr}
    \toprule
    Statistic & Observed & Null mean $\pm$ SD & $p$ \\
    \midrule
    Block ratio       & 1.581 & $1.001 \pm 0.025$ & $< 0.001$ \\
    Cross-era $\Wass$ & 0.752 & $0.618 \pm 0.007$ & $< 0.001$ \\
    Within-BHPS       & 0.486 & $0.617 \pm 0.014$ & $< 0.001$ \\
    Within-USoc       & 0.466 & $0.618 \pm 0.019$ & $< 0.001$ \\
    Consecutive mean  & 0.467 & $0.618 \pm 0.027$ & $< 0.001$ \\
    Mantel $r$        & 0.768 & $0.001 \pm 0.048$ & $< 0.001$ \\
    \bottomrule
  \end{tabular}
\end{table}

\textbf{Pool-draw full} (50 permutations, 32 diagrams each).
The block structure cannot be reproduced by random year-ignorant subsampling:
\begin{itemize}
  \item Block ratio: observed 1.581 vs null $1.007 \pm 0.095$ ($p < 0.001$)
  \item Cross-era $\Wass$: observed 0.752 vs null $0.481 \pm 0.042$
    ($p < 0.001$)
\end{itemize}
However, consecutive-year smoothness is \emph{not} significant (observed 0.467
vs null $0.476 \pm 0.026$, $p = 0.34$): adjacent years' diagrams are no more
similar than random subsets.
The dominant signal is the era-level block, not fine-grained year-to-year
evolution.

\begin{figure}[htbp]
  \centering
  \includegraphics[width=\linewidth]{fig5_mds_embedding}
  \caption{MDS embedding of the Wasserstein distance matrix, coloured by survey
    era (BHPS circles, USoc squares), with recession years marked by red outlines.
    The clear left-right era separation contrasts with the absence of a distinct
    GFC cluster within the USoc era.}
  \label{fig:mds}
\end{figure}

\subsubsection{$\HH{1}$ Block Structure}

Loop topology ($\HH{1}$) shows minimal era separation: cross-era $\Wass$ is only
13\% above within-era means.
The loop structure of the trajectory space is stable across the survey
transition; the design confound operates almost entirely through connected
components, not loops.

\subsection{Macroeconomic Correlations}
\label{sec:macro}

Having established that the dominant cross-era signal is survey design, we
restrict the most informative macro analysis to within-era correlations.
Full-panel analyses at four robustness levels supplement this.

\subsubsection{Full-Panel Correlations}

Seventeen of 42 raw Pearson correlations survive BH-FDR correction.
The strongest (Table~\ref{tab:macro_full}) are uniformly labour-market
\emph{composition} variables, not \emph{performance} variables.

\begin{table}[htbp]
  \centering
  \caption{Strongest raw correlations surviving BH-FDR ($\alpha = 0.05$).}
  \label{tab:macro_full}
  \begin{tabular}{llrr}
    \toprule
    Macro indicator & Topological measure & $r$ & $p$ \\
    \midrule
    Self-employment share & $\AUC(\beta_0)$ & $+0.872$ & $< 0.0001$ \\
    Self-employment share & $\TP(\HH{0})$   & $+0.846$ & $< 0.0001$ \\
    Self-employment share & $\PE(\HH{0})$   & $-0.792$ & $< 0.0001$ \\
    Gini                  & $\PE(\HH{0})$   & $+0.740$ & $< 0.0001$ \\
    Employment rate       & $\TP(\HH{0})$   & $+0.698$ & $< 0.0001$ \\
    Employment rate       & $\AUC(\beta_0)$ & $+0.638$ & $0.0001$ \\
    Gini                  & $\AUC(\beta_0)$ & $-0.590$ & $0.0004$ \\
    Part-time share       & $\TP(\HH{0})$   & $+0.559$ & $0.0009$ \\
    \bottomrule
  \end{tabular}
\end{table}

GDP growth and inflation show no significant association with any topological
measure after correction.

\subsubsection{Robustness Hierarchy}

The four specifications reveal a clear hierarchy (Table~\ref{tab:robustness}).

\begin{table}[htbp]
  \centering
  \caption{Significant correlations by specification (BH-FDR at $\alpha = 0.05$).}
  \label{tab:robustness}
  \begin{tabular}{lrl}
    \toprule
    Specification & Surviving / Total & Interpretation \\
    \midrule
    Raw             & 17 / 42 & Many associations (includes shared trends) \\
    Detrended       & 10 / 35 & Most survive; removes linear time trends \\
    Within-BHPS     &  4 / 35 & Pure within-survey, constant design \\
    First-differenced &  0 / 35 & No year-on-year associations \\
    \bottomrule
  \end{tabular}
\end{table}

The critical result: year-on-year changes in macroeconomic indicators do not
predict year-on-year changes in topology.
Topology tracks slow structural shifts in \emph{who works and how}, not annual
business-cycle fluctuations.

\subsubsection{Within-BHPS Correlations}

Within the BHPS era alone---where survey design is constant---the following
associations survive (Table~\ref{tab:within_bhps}).

\begin{table}[htbp]
  \centering
  \caption{Significant correlations within the BHPS era (BH-FDR, $\alpha = 0.05$).}
  \label{tab:within_bhps}
  \begin{tabular}{llrr}
    \toprule
    Macro indicator & Topological measure & $r$ & $p$ \\
    \midrule
    Employment rate & $\AUC(\beta_0)$ & $+0.823$ & $< 0.001$ \\
    Unemployment    & $\AUC(\beta_0)$ & $-0.807$ & $< 0.001$ \\
    Part-time share & $\AUC(\beta_0)$ & $+0.795$ & $< 0.001$ \\
    Employment rate & $\TP(\HH{0})$   & $+0.763$ & $< 0.001$ \\
    \bottomrule
  \end{tabular}
\end{table}

These are strong associations within a constant survey design: as the employment
rate rises and unemployment falls during the 1990s--2000s UK expansion, the
trajectory space becomes more topologically complex (more fine-scale
fragmentation).
A tighter labour market with rising part-time and self-employment produces a more
heterogeneous population of career trajectories, a broader menu of employment
configurations that manifests as increased topological complexity.

\subsubsection{Detrended and Within-USoc}

After detrending, unemployment flips sign (raw $r = -0.45 \to$ detrended
$r = +0.69$), which makes sense: cyclically, unemployment spikes coincide with
transient fine-scale fragmentation.

Within USoc, self-employment ($r = +0.71$) and inflation ($r = -0.64$) are the
only significant correlates, a narrower set consistent with the USoc era's
topological homogeneity observed in the Mantel test.

\subsubsection{Lagged Analysis}

No significant lag structure emerges for GDP growth at any lead/lag ($\pm 3$
years).
Unemployment shows marginal strengthening when topology leads ($r = -0.44$ at
lag $+3$), but this is likely a level-shift artefact rather than causal dynamics.
We do not claim that topology functions as a leading indicator.

%%%%%%%%%%%%%%%%%%%%%%%%%%%%%%%%%%%%%%%%%%%%%%%%%%%
\section{Discussion}
\label{sec:discussion}
%%%%%%%%%%%%%%%%%%%%%%%%%%%%%%%%%%%%%%%%%%%%%%%%%%%

\subsection{The Cautionary Finding}

Zigzag persistence, and any time-varying topological method, does not on its own
distinguish substantive distributional change from survey-design artefacts.
The na\"ive zigzag analysis (\S\ref{sec:naive}) produced an appealing story about
the 2008 financial crisis creating a permanent topological scar in the UK
trajectory space.
The spanning-individual decomposition (\S\ref{sec:decomp}) showed this story to
be mostly wrong.

This vulnerability is not specific to zigzag persistence.
Any distributional comparison (Wasserstein distances, kernel MMD tests, classifier
two-sample tests) would detect the BHPS--USoc shift.
What the topological framework gives is \emph{structured} decomposition: the
block structure in the Wasserstein matrix, the era-specific zigzag, the within-era
Mantel test, and the loop-stability finding each supply a distinct piece of
evidence.
Together, they allow confident source identification.

For applied time-varying TDA, the implication is that whenever point-cloud
populations change for non-substantive reasons (survey redesigns, missingness
patterns, covariate shift), topological change should be assumed partly artefactual
until proven otherwise.
The spanning-individual decomposition provides one route to this proof.

\subsection{What Survives: Secular Flexibilisation}

Having removed the survey confound, the remaining signal is worth attention.
Within the BHPS era (1991--2008), total $\HH{0}$ persistence increases at
$+0.31$/year ($p < 0.0001$), and the CUSUM changepoint at 2006 places the
inflection before the GFC.
The within-BHPS correlation with employment rate ($r = +0.82$), unemployment
($r = -0.81$), and part-time share ($r = +0.80$) is strong and robust.

This pattern aligns with the flexibilisation of UK labour markets during
1991--2008: rising employment rates, growth of part-time and temporary work,
expansion of self-employment, and the spread of non-standard employment
arrangements \citep{rubery2016labour, grimshaw2017reducing}.
As the labour market tightened and diversified, the space of career trajectories
grew more complex---not because individual lives became more chaotic, but because
the range of populated employment configurations expanded.
More trajectory types gained members; the point cloud grew more geometrically
diverse.

First-differenced correlations yield nothing, which tells us this is a slow
structural process, not a business-cycle phenomenon.
Topological complexity tracks \emph{levels} of employment diversity, not
\emph{changes} in GDP.
This is the kind of long-run structural shift that the sociological literature on
flexibilisation describes, now given a topological characterisation.

\subsection{The Wasserstein Distance as a Survey-Comparability Diagnostic}

The $32 \times 32$ Wasserstein distance matrix has methodological value beyond
this paper.
Two surveys that produce similar persistence diagrams yield low $\Wass$; a high
cross-era $\Wass$, especially one exceeding within-era distances, flags a design
discontinuity.
Here the block ratio (1.58) and perfect hierarchical-clustering separation at
$k = 2$ provide a clear diagnostic.

This suggests a broader application: for any linked longitudinal study, computing
pairwise Wasserstein distances between annual topological summaries could work as
a routine comparability check.
A topological distance matrix (as in Figure~\ref{fig:heatmap}) would flag
design-related breaks, possibly more sensitively than standard balance tests on
marginal distributions, because persistence diagrams capture geometric structure
that marginals miss.

\subsection{$\HH{1}$ Stability}

The minimal cross-era $\HH{1}$ difference (13\% elevation vs 55--60\% for
$\HH{0}$) deserves attention.
Loop structures in the trajectory space, corresponding to cyclical career
patterns (cycling between employment and unemployment, or between income bands),
appear invariant to both survey design and macroeconomic conditions.
The \emph{qualitative forms} of cyclical career dynamics are stable over three
decades, even as the \emph{number and distribution} of trajectory clusters
($\HH{0}$) change.
Whether this extends to higher homology is an open question.

\subsection{Limitations}

\textbf{Zigzag landmark sensitivity.}
The 4-bar structure at $\varepsilon = 1.0$ requires $L \geq 100$ landmarks; at
$L = 75$, only one bar survives.
The sensitivity grid shows this is $\varepsilon$-dependent, not a resolution
issue, but formal bootstrap stability analysis on bar lifetimes would strengthen
the finding.

\textbf{Frozen PCA.}
The common coordinate system depends on PCA loadings fitted to the full (BHPS +
USoc) sample.
If the factor structure genuinely differs between surveys, frozen loadings may
distort.
Robustness checks with era-specific PCA would address this.

\textbf{Causal inference.}
With 32 time points, and only 18 within BHPS, statistical power for dynamic
specifications is limited.
The within-BHPS correlations are descriptive associations, not causal claims.

\textbf{State-space coarseness.}
The nine-state employment $\times$ income classification discards within-state
heterogeneity.
A finer classification (adding sector or occupation) might reveal richer
topological structure, but at the cost of sparser $n$-grams.

\textbf{UK-specific.}
The BHPS--USoc transition is specific to the UK.
The methodological lesson (survey confounds in time-varying TDA) generalises; the
substantive findings may not.

\subsection{Connections to the Research Programme}

Paper~1 \citep{dorman2026a} establishes static topological structure in the
trajectory space.
Paper~2 \citep{dorman2026b} applies Mapper to reveal its interior graph structure.
This paper adds the temporal dimension and the cautionary finding about survey
design.

The Wasserstein distance matrix opens paths for Papers~5 (cross-national
comparison) and~6 (intergenerational analysis): pairwise $\Wass$ between
countries' annual diagrams would provide a topological metric of welfare-state
similarity, and parent-child $\Wass$ comparisons could measure intergenerational
topological mobility.
In each case, the diagnostics developed here (spanning populations, block-structure
tests, pool-draw nulls) would be needed for separating design from substance.

%%%%%%%%%%%%%%%%%%%%%%%%%%%%%%%%%%%%%%%%%%%%%%%%%%%
\section{Conclusion}
%%%%%%%%%%%%%%%%%%%%%%%%%%%%%%%%%%%%%%%%%%%%%%%%%%%

Zigzag persistent homology applied to 32 years of UK employment-income
trajectories reveals that what appeared to be a topological signature of the 2008
financial crisis is predominantly a survey-design effect: Understanding Society's
broader sampling frame captures career types absent from the British Household
Panel Study.
Establishing this required era-specific robustness tests, spanning-individual
analysis, and Wasserstein null models, confirming that time-varying topological
methods in the social sciences need the same attention to survey methodology as
any other longitudinal analysis.

The harder-won finding is a genuine secular increase in topological complexity
within the BHPS era (1991--2008), driven by labour market composition rather than
business-cycle dynamics.
Topology tracks the slow structural diversification of UK careers, a quantitative
expression of flexibilisation that was underway before the financial crisis and
already visible in the geometry of the trajectory space.

The method works.
It detects real distributional change.
The challenge is that not all distributional change is substantively interesting,
and decomposing the signal requires tools beyond topology alone.

\begin{figure}[htbp]
  \centering
  \includegraphics[width=\linewidth]{fig6_zigzag_schematic}
  \caption{Zigzag tower (diamond construction) over annual snapshot sub-clouds.
    Each annual point cloud $P_t$ is linked to its neighbours via the union
    complexes $P_t \cup P_{t+1}$, enabling topological features to be tracked
    across time without requiring a common simplex filtration.}
  \label{fig:schematic}
\end{figure}
